\documentclass[prb,amsmath,amssymb,onecolumn,nopacs,floatfix,longbibliography]{revtex4-1}
\usepackage{graphicx}
\usepackage[english]{babel}
\usepackage{times}
\usepackage{dcolumn}
\usepackage[usenames,dvipsnames]{color}
\usepackage{units}
\usepackage{color}

\begin{document}

\title{Evaluation of integrals in transport calculations within the DMFT}

\author{Rok \v{Z}itko}

\date{\today}

\begin{abstract}
\end{abstract}

%\pacs{72.10.Fk, 72.15.Qm}

% 72.10.Fk Scattering by point defects, dislocations, surfaces, and
% other imperfections (including Kondo effect)
% 72.15.Qm Scattering mechanisms and Kondo effect (see also 75.20.Hr
% Local moments in compounds and alloys; Kondo effect, valence
% fluctuations, heavy fermions in magnetic properties and materials)
% 72.25.Dc Spin polarized transport in semiconductors

\maketitle

\newcommand{\vc}[1]{{\mathbf{#1}}}
\newcommand{\vck}{\vc{k}}
\newcommand{\braket}[2]{\langle#1|#2\rangle}
\newcommand{\expv}[1]{\langle #1 \rangle}
\newcommand{\corr}[1]{\langle\langle #1 \rangle\rangle}
\newcommand{\ket}[1]{| #1 \rangle}
\newcommand{\Tr}{\mathrm{Tr}}
\newcommand{\kor}[1]{\langle\langle #1 \rangle\rangle}
\newcommand{\degg}{^\circ}
\renewcommand{\Im}{\mathrm{Im}}
\renewcommand{\Re}{\mathrm{Re}}
\newcommand{\dtN}{{\dot N}}
\newcommand{\dtQ}{{\dot Q}}
\newcommand{\sgn}{\mathrm{sgn}}
\newcommand{\beq}[1]{\begin{equation} #1 \end{equation}}
\newcommand{\one}{\mathbf{1}}


\section{Introduction}

The behavior of single-particle excitations in a material can be
theoretically described using the formalism of Green's functions. We
define the spectral function $\rho$ and the Green's function $G$ as
\beq{
\begin{split}
\rho(\omega,\vc{k}) &= -\frac{1}{\pi} \Im G(\omega,\vc{k}), \\
G(\omega,\vc{k}) &= \frac{1}{\omega+\mu-\epsilon_\vc{k}-\Sigma(\omega,\vc{k})},
\end{split}
}
%
where $\epsilon_\vc{k}$ is the energy of an electron with momentum
$\vc{k}$, $\omega$ is frequency, $\mu$ is the chemical potential, and
$\Sigma(\omega,\vc{k})$ is called the self-energy function which
describes the effects of interactions.
$\Sigma(\omega,\vc{k})$ is a complex function such that
$\Im\Sigma(\omega,\vc{k})<0$ for all arguments, which is a consequence
of causality and guarantees that the spectral function
$\rho(\omega,\vc{k})$ is a positive quantity. 

The conductivity of a material is related to the current-current
correlations (Kubo's formula). The current, in turn, is proportional
to the velocity of electrons $v(\vc{k})=d\epsilon_\vc{k}/dk$ and their
distribution function $A(\vc{k})$ (the fraction of electrons with
momentum $\vc{k}$). The correlators are quadratic in $v(\vc{k})$ and
$A(\vc{k})$, and we need to sum over contributions from all $\vc{k}$.
In some cases, the $\vc{k}$ dependence enters only through the
dependence of $\epsilon_\vc{k}$, and the $\vc{k}$ sum can be replaced
by an integral over electron energy $\epsilon=\epsilon_\vc{k}$.

In the limit of infinite dimensions, the self-energy becomes purely
local, $\Sigma=\Sigma(\omega)$, and there are no vertex corrections.
The conductivity can then be evaluated from a simple double
integration over electron energy $\epsilon$ and frequency $\omega$
(bubble formula). This holds for other transport properties as well.
For metals, the integrals become difficult to evaluate in the limit of
low temperatures because of sharply peaked (almost singular)
contributions close to the Fermi level, because for a Fermi-liquid
system the self-energy $\Sigma$ behaves as
%
\beq{
\Sigma(\omega) \approx \Re \Sigma(0) + \omega \left( 1- \frac{1}{Z}
\right) - i \frac{1}{Z^2} [ (\pi k_B T)^2 + \omega^2 ],
}
%
thus the broadening ($\Im\Sigma$) becomes small in the $\omega \to 0$,
$T \to 0$ limit. Here $0<Z\leq 1$ is the quasiparticle residue, $k_B$
is the Boltzmann constant, and $T$ is the temperature. More generally,
such intergrals are difficult to evaluate whenever $\Im \Sigma$ becomes
small.

Our goal is to perform the $\epsilon$ integration in the bubble
formula analytically and reduce the problem to a simpler numerical
quadrature over $\omega$ only. We need an accurate and precise result
in short time. The code needs to be robust with respect to
$\Sigma(\omega)$ because the self-energy can be affected with
numerical noise and is not necessarily very smooth.

In the following $k_B=\hbar=1$.


\section{Problem definition}

\newcommand{\de}{\mathrm{d}\epsilon}
\newcommand{\dom}{\mathrm{d}\omega}

We consider the (dimensionless) integrals
%
\beq{
I_{mno} = \frac{1}{D} \int_{-\infty}^\infty \de\, \Phi_m(\epsilon)
\int_{-\infty}^{+\infty} \dom \left( -\frac{\partial f(\omega)}{\partial
\omega} \right) \left[ D \rho(\omega,\epsilon) \right]^n 
\left( \frac{\omega}{D} \right)^o,
}
%
where $\Phi_m(\epsilon)$ is one of the (dimensionless) functions
%
\beq{
\begin{split}
\Phi_1(\epsilon) &= 1 \times b(\epsilon), \\
\Phi_2(\epsilon) &= \epsilon/D \times b(\epsilon), \\
\Phi_3(\epsilon) &= [1-(\epsilon/D)^2]^{1/2} \times b(\epsilon), \\
\Phi_4(\epsilon) &= (\epsilon/D) [1-(\epsilon/D)^2]^{1/2} \times b(\epsilon), \\
\Phi_5(\epsilon) &= [1-(\epsilon/D)^2]^{3/2} \times b(\epsilon), \\
\Phi_6(\epsilon) &= (\epsilon/D) [1-(\epsilon/D)^2]^{3/2} \times b(\epsilon), \\
\Phi_7(\epsilon) &= \exp[-2 (\epsilon/\sigma)^2], \\
\Phi_8(\epsilon) &= (\epsilon/\sigma) \exp[-2(\epsilon/\sigma)^2].
\end{split}
}
%
$b(\epsilon)$ is a box function which is equal to 1 in the interval
$[-D:D]$ and zero otherwise. For $m=1,\ldots,6$ the $\epsilon$
integral thus extends from $-D$ to $D$, for $m=7,8$ it extends over
the full real axis. For flat band ($m=1,2$) and the Bethe lattice
($m=3-7$), $D$ has the meaning of the band half-width, for
hypercubic lattice ($m=7,8$) it is just a rescale factor. By default,
$D=1$ and $\sigma=1$, which are commonly used conventions. The
expressions with the factor of $\epsilon$ are used in the expressions
for the Hall transport \cite{arsenault2013}.

$f(\omega)=1/(1+e^{\omega/T})$ is the Fermi function at temperature
$T$. We use the convention that the argument $\omega=0$ corresponds to
the chemical potential in $G$, $\rho$ and $f$. The chemical potential
$\mu$ only appears as the shift in the denominator of $G$.

$n$ and $o$ are some integers. $n=2,3$, $o=0,1,2$ are of particular
interest. For conductivity, $n=2$, $o=0$. For thermopower, $n=2$,
$o=1$. For thermal transport, $n=2$, $o=2$. For Hall conductivity,
$n=3$, $o=1$. For Hall thermoconductivity, $n=3$, $o=2$. (The code
does not deal with the prefactors to $I_{mno}$ to find these physical
quantities, it just computes the dimensionless integral $I_{mno}$ and
it is up to the user to decide what to do with the resulting number.
In particular, it is to be noted that $\rho$ is defined here as the
spectral function including the $\pi$ factor.)

The input are the indexes $m$, $n$ and $o$, temperature $T$, chemical
potential $\mu$, and tabulated $\Sigma(\omega)$, i.e., a file
containing columns $\{\omega,\Re\Sigma(\omega),\Im\Sigma(\omega)\}$.

The output is $I_{mno}$ (real quantity) and optionally an error
estimate.


\section{Implementation}

The $\epsilon$ integrals
%
\beq{
J_{mn}(\Omega) = \frac{1}{D} \int_{-\infty}^\infty \de\, \Phi_m(\epsilon)
\left[ -\frac{1}{\pi} \Im \frac{D}{\Omega-\epsilon} \right]^{n}
}
%
are evaluated analytically in Mathematica for all $m$ and $n=2,3$.
These twelve (dimensionless) functions take $\Omega$ as input and
return a real quantity as output. It may be assumed that $\Im \Omega >
0$ (and the code needs to assert that this is indeed the case). In
practice, we will always take $\Omega=\omega-\Sigma(\omega)$ as
argument.

The resulting expressions are to be converted to C code and saved to a
file which will be included from the main program written in C or C++.
This step of the code generation has to be fully automated (no
copy/pasting!) and performed, for example, through evaluating a
Mathematica notebook or script.

In the main program, $\Sigma(\omega)$ is read from a file and
interpolated; the user can choose between linear and higher-order
(cubic spline, Akima spline) interpolation. GNU scientific library (GSL)
interpolation object are used for this purpose.

\newcommand{\mn}{\mathrm{min}}
\newcommand{\mx}{\mathrm{max}}
\newcommand{\dx}{\mathrm{d}x}

The $\omega$ integration
%
\beq{
I_{mno} = \int_{-\infty}^{+\infty} \dom \left(
-\frac{\partial f(\omega)}{\partial \omega} \right)
J_{mn}[\omega-\Sigma(\omega)]
(\omega/D)^o
}
%
is then performed. The tabulated $\Sigma$ defines the integral
boundaries ($\omega_\mn$ and $\omega_\mx$) to be used instead of
($-\infty$, $+\infty$). In other words, $\Im\Sigma$ is assumed to be
zero outside the interpolation interval.

The derivative of the Fermi function can be written as
%
\beq{
-\frac{\partial f(\omega)}{\partial\omega} = \frac{1}
{2T(1+\cosh x)}
}
%
where $x=\omega/T$. Because of this weight function, the main
contribution to the integral comes from the interval $x \in [-5:5]$.
The idea is thus to rewrite the integral $I_{mn}$ in terms of the
variable $x$,
%
\beq{
I_{mno} = \int \dx\, g(x) J_{mn}[T x - \Sigma(T x)] \left(
\frac{T x}{D} \right)^o,
}
%
where $g(x)=1/(2+2\cosh(x))$. The suitable upper limits to $x$
integration have to be recalculated from $\omega_\mn$ and
$\omega_\mx$, but can be constrained within some sensible limits
(e.g., -40 and 40) because $-\partial f/\partial \omega$ is
exponentially decreasing and beyond 40 it is smaller than numerical
round-off errors.

The integration is to be performed with the GSL quadrature routines
with adaptive grid algoritms. It would be of interest to separately
perform the integration for $x>0$ and $x<0$, with an option to display
the partial results (if verbosity is increased using {\tt -v}, see
below).


\section{Invocation}

The executable is named {\tt bubble}. The syntax is:

{\tt bubble [-v] [-i <i>] [-p <p>] [-D <D>] [-s <sigma>] 
m n o T mu Sigmafile}

The switch {\tt -v} enables verbose output (e.g. error estimates). By
default, only the final result $I_{mno}$ is given as the output.

{\tt i} controls the order of interpolation: $i=1$ is linear
interpolation, $i=2$ is quadratic interpolation, $i=3$ is Akima spline
interpolation.

{\tt p} controls the precision of the numerical integration.

The defaults are $D=1$, $\sigma=1$. They can be changed to
{\tt D} and {\tt sigma}.



\section{Testing}

The code can be tested using the limit of constant relaxation time,
$\Sigma(\omega)=-i\Gamma$ with some constant $\Gamma>0$. Then
$I_{mno}$ can be evaluated analytically for different $T$ and the
results compared to the output of {\tt bubble}. A suitable
input file then consists simply of two lines:

{\tt omega\_min 0 -Gamma}

{\tt omega\_max 0 -Gamma}

Testing also has to be automated and rerunnable (to be useful as
a regression test).

Additionally, we want to test using real-life examples of
$\Sigma(\omega)$ taken from DMFT(NRG) calculations for the Hubbard
model, and compare against the known results for the transport quantities.


\section{Concluding remarks}

In the future we might be interested in additional functions
$\Phi_m(\epsilon)$ corresponding to different lattices (e.g. square,
cubic).


\nocite{bevilacqua2013}

\bibliography{bubble}


\end{document}


 For non-interacting electron, the conductivity is
hence proportional to
%
\beq{
\label{bubble}
\sum_\vc{k} v_\vc{k}^2 A(\vc{k})^2 = \int \mathrm{d}\epsilon \sum_\vc{k}
\delta(\epsilon-\epsilon_k) v(\epsilon)^2 A(\epsilon)^2
= \int \mathrm{d}\epsilon\, r(\epsilon) v^2(\epsilon)^2 A(\epsilon)^2,
}
%
the assumption here being that $v(\vc{k})$ and $A(\vc{k})$ depend on
$\vc{k}$ only through the energy $\epsilon_\vc{k}$. We have introduced
the density of states functions $r(\epsilon)=\sum_\vc{k}
\delta(\epsilon-\epsilon_\vc{k})$. Under certain conditions, a similar
expression as Eq.~\eqref{bubble} can be written for interacting systems
in terms of the spectral function $\rho$ in place of the distribution
function $A$.

